\newglossaryentry{prompt}
{
    name=prompt,
    description={Instrucción o pregunta en lenguaje natural que se proporciona a un modelo de lenguaje para obtener una respuesta específica. Es la interfaz principal de comunicación entre el usuario y el LLM}
}

\newglossaryentry{chunk}
{
    name=chunk,
    description={Fragmento o segmento de texto de tamaño limitado utilizado para dividir documentos largos. Los chunks permiten procesar información de manera eficiente en modelos con limitaciones de contexto}
}

\newglossaryentry{embedding}
{
    name=embedding,
    description={Representación vectorial numérica de texto que captura su significado semántico. Permite a los modelos comparar y buscar información basándose en similitud conceptual}
}

\newglossaryentry{token}
{
    name=token,
    description={Unidad mínima de texto procesada por un modelo de lenguaje. Puede ser una palabra completa, parte de una palabra, o un carácter, dependiendo del modelo de tokenización utilizado}
}

\newglossaryentry{contexto}
{
    name=ventana de contexto,
    description={Cantidad máxima de tokens que un modelo de lenguaje puede procesar simultáneamente. Define el límite de información que el modelo puede considerar al generar una respuesta}
}

\newglossaryentry{inferencia}
{
    name=inferencia,
    description={Proceso mediante el cual un modelo entrenado genera respuestas o predicciones basándose en nuevos datos de entrada, sin modificar sus parámetros internos}
}

\newglossaryentry{fine-tuning}
{
    name=fine-tuning,
    description={Proceso de reentrenamiento de un modelo preentrenado con datos específicos de un dominio para mejorar su desempeño en tareas particulares}
}

\newglossaryentry{alucinacion}
{
    name=alucinación,
    description={Fenómeno donde un modelo de lenguaje genera información incorrecta, inventada o sin fundamento, presentándola con aparente confianza como si fuera verdadera}
}

\newglossaryentry{temperatura}
{
    name=temperatura,
    description={Parámetro que controla la aleatoriedad en las respuestas del modelo. Valores bajos producen respuestas más determinísticas, mientras que valores altos generan mayor creatividad y variabilidad}
}

\newglossaryentry{api}
{
    name=API,
    description={Application Programming Interface. Conjunto de definiciones y protocolos que permite la comunicación entre diferentes componentes de software}
}

\newglossaryentry{vectordb}
{
    name=base de datos vectorial,
    description={Sistema de almacenamiento optimizado para guardar y buscar embeddings de manera eficiente mediante cálculos de similitud vectorial}
}

\newglossaryentry{transcripcion}
{
    name=transcripción,
    description={Proceso de convertir audio hablado en texto escrito mediante sistemas de reconocimiento automático de voz (ASR)}
}

\newglossaryentry{accuracy}
{
    name=accuracy,
    description={Métrica que mide la proporción de predicciones correctas sobre el total de predicciones realizadas por un modelo}
}

\newglossaryentry{frontend}
{
    name=frontend,
    description={Parte de una aplicación web con la que los usuarios interactúan directamente. Incluye la interfaz visual y la lógica de presentación}
}

\newglossaryentry{backend}
{
    name=backend,
    description={Parte del servidor de una aplicación que maneja la lógica de negocio, procesamiento de datos, y comunicación con bases de datos}
}

\newglossaryentry{react}
{
    name=React,
    description={Biblioteca de JavaScript desarrollada por Facebook para construir interfaces de usuario mediante componentes reutilizables}
}

\newglossaryentry{nodejs}
{
    name=Node.js,
    description={Entorno de ejecución de JavaScript del lado del servidor, construido sobre el motor V8 de Chrome}
}

\newglossaryentry{OpenFing}
{
    name=OpenFing,
    description={Plataforma educativa de la Facultad de Ingeniería de la Universidad de la República que ofrece contenido académico abierto, principalmente videos de clases}
}

% Acrónimos
\newacronym{llm}{LLM}{Large Language Model}
\newacronym{rag}{RAG}{Retrieval-Augmented Generation}
\newacronym{nlp}{NLP}{Natural Language Processing}
\newacronym{asr}{ASR}{Automatic Speech Recognition}
\newacronym{ia}{IA}{Inteligencia Artificial}
\newacronym{json}{JSON}{JavaScript Object Notation}
\newacronym{http}{HTTP}{Hypertext Transfer Protocol}
\newacronym{url}{URL}{Uniform Resource Locator}
\newacronym{csv}{CSV}{Comma-Separated Values}
\newacronym{ui}{UI}{User Interface}
\newacronym{ux}{UX}{User Experience}


